% --- INICIO DE PLANTILLA PERSONALIZADA ---
\documentclass[oneside,11pt]{Latex/Classes/PhDthesisPSnPDF}

% --- Paquetes Originales ---
\usepackage{blindtext}
\usepackage{actuarialangle} 
\usepackage{tabularx}
\usepackage{float}
\usepackage{multirow, array}
\usepackage[usenames,dvipsnames,svgnames,table]{xcolor}
\usepackage{longtable}
\usepackage{latexsym,amsmath,amssymb,amsfonts}
\usepackage{enumitem}
\usepackage{xparse}
\usepackage{booktabs}
\usepackage[round, sort, numbers]{natbib}  
\usepackage{adjustbox}
\usepackage[utf8]{inputenc}
\usepackage{caption}
\usepackage{graphicx}
\usepackage{hyperref}
\usepackage{cleveref}

% Definición de colores
\definecolor{miblue}{RGB}{68, 114, 196}

% --- CAMBIO CRÍTICO 1: Usar \input para macros en el preámbulo ---
% \include da problemas antes del begin document. \input es seguro.
\input{Latex/Macros/MacroFile1}

% --- Definiciones de seguridad (por si la clase falla al cargar alguna variable) ---
\providecommand{\facultad}[1]{\def\lafacultad{#1}}
\providecommand{\escudofacultad}[1]{\def\elescudo{#1}}
\providecommand{\degree}[1]{\def\elgrado{#1}}
\providecommand{\director}[1]{\def\eldirector{#1}}
\providecommand{\lugar}[1]{\def\ellugar{#1}}
\providecommand{\degreedate}[1]{\def\lafecha{#1}}

% --- Configuración para bloques de código de R (Pandoc) ---
\usepackage{color}
\usepackage{fancyvrb}
\newcommand{\VerbBar}{|}
\newcommand{\VERB}{\Verb[commandchars=\\\{\}]}
\DefineVerbatimEnvironment{Highlighting}{Verbatim}{commandchars=\\\{\}}
% Add ',fontsize=\small' for more characters per line
\usepackage{framed}
\definecolor{shadecolor}{RGB}{248,248,248}
\newenvironment{Shaded}{\begin{snugshade}}{\end{snugshade}}
\newcommand{\AlertTok}[1]{\textcolor[rgb]{0.94,0.16,0.16}{#1}}
\newcommand{\AnnotationTok}[1]{\textcolor[rgb]{0.56,0.35,0.01}{\textbf{\textit{#1}}}}
\newcommand{\AttributeTok}[1]{\textcolor[rgb]{0.13,0.29,0.53}{#1}}
\newcommand{\BaseNTok}[1]{\textcolor[rgb]{0.00,0.00,0.81}{#1}}
\newcommand{\BuiltInTok}[1]{#1}
\newcommand{\CharTok}[1]{\textcolor[rgb]{0.31,0.60,0.02}{#1}}
\newcommand{\CommentTok}[1]{\textcolor[rgb]{0.56,0.35,0.01}{\textit{#1}}}
\newcommand{\CommentVarTok}[1]{\textcolor[rgb]{0.56,0.35,0.01}{\textbf{\textit{#1}}}}
\newcommand{\ConstantTok}[1]{\textcolor[rgb]{0.56,0.35,0.01}{#1}}
\newcommand{\ControlFlowTok}[1]{\textcolor[rgb]{0.13,0.29,0.53}{\textbf{#1}}}
\newcommand{\DataTypeTok}[1]{\textcolor[rgb]{0.13,0.29,0.53}{#1}}
\newcommand{\DecValTok}[1]{\textcolor[rgb]{0.00,0.00,0.81}{#1}}
\newcommand{\DocumentationTok}[1]{\textcolor[rgb]{0.56,0.35,0.01}{\textbf{\textit{#1}}}}
\newcommand{\ErrorTok}[1]{\textcolor[rgb]{0.64,0.00,0.00}{\textbf{#1}}}
\newcommand{\ExtensionTok}[1]{#1}
\newcommand{\FloatTok}[1]{\textcolor[rgb]{0.00,0.00,0.81}{#1}}
\newcommand{\FunctionTok}[1]{\textcolor[rgb]{0.13,0.29,0.53}{\textbf{#1}}}
\newcommand{\ImportTok}[1]{#1}
\newcommand{\InformationTok}[1]{\textcolor[rgb]{0.56,0.35,0.01}{\textbf{\textit{#1}}}}
\newcommand{\KeywordTok}[1]{\textcolor[rgb]{0.13,0.29,0.53}{\textbf{#1}}}
\newcommand{\NormalTok}[1]{#1}
\newcommand{\OperatorTok}[1]{\textcolor[rgb]{0.81,0.36,0.00}{\textbf{#1}}}
\newcommand{\OtherTok}[1]{\textcolor[rgb]{0.56,0.35,0.01}{#1}}
\newcommand{\PreprocessorTok}[1]{\textcolor[rgb]{0.56,0.35,0.01}{\textit{#1}}}
\newcommand{\RegionMarkerTok}[1]{#1}
\newcommand{\SpecialCharTok}[1]{\textcolor[rgb]{0.81,0.36,0.00}{\textbf{#1}}}
\newcommand{\SpecialStringTok}[1]{\textcolor[rgb]{0.31,0.60,0.02}{#1}}
\newcommand{\StringTok}[1]{\textcolor[rgb]{0.31,0.60,0.02}{#1}}
\newcommand{\VariableTok}[1]{\textcolor[rgb]{0.00,0.00,0.00}{#1}}
\newcommand{\VerbatimStringTok}[1]{\textcolor[rgb]{0.31,0.60,0.02}{#1}}
\newcommand{\WarningTok}[1]{\textcolor[rgb]{0.56,0.35,0.01}{\textbf{\textit{#1}}}}

%%%%%%%%%%%%%%%%%%%%%%%%%%%%%%%%%%%%%%%%%%%%%%%%%%%%%%%%%%%%%%%%%%%%%%%%%%%%%%%%
%                         DATOS (Conectados a RMarkdown)                       %
%%%%%%%%%%%%%%%%%%%%%%%%%%%%%%%%%%%%%%%%%%%%%%%%%%%%%%%%%%%%%%%%%%%%%%%%%%%%%%%%
\title{Título del Trabajo Final}
\author{Rehide De Pablos \\ Moisés Perez \\ Diego Rosabrussin}

% Lógica para Facultad: Si la defines en el YAML la usa, si no, usa la default

\facultad{Facultad de Ciencias Económicas y Sociales \\ Escuela de
Estadistica y Ciencias Actuariales}
  \escudofacultad{Latex/Classes/Escudos/faces}

\degree{Computación I}
\director{Jesus Ochoa \\ Oliver Triveño}
\degreedate{Febrero 2026} 
\lugar{Caracas}

\portadatrue 

% Metadatos PDF
\keywords{}
\subject{}

% --- CORRECCIÓN PARA PANDOC NUEVO (Imágenes) ---
\makeatletter
\providecommand{\pandocbounded}[1]{#1}
\makeatother


%%%%%%%%%%%%%%%%%%%%%%%%%%%%%%%%%%%%%%%%%%%%%%%%%%%%%
%                   DOCUMENTO                       %
%%%%%%%%%%%%%%%%%%%%%%%%%%%%%%%%%%%%%%%%%%%%%%%%%%%%%
\begin{document}

\maketitle

%%%%%%%%%%%%%%%%%%%%%%%%%%%%%%%%%%%%%%%%%%%%%%%%%%%%%
%                  PRÓLOGO                          %
%%%%%%%%%%%%%%%%%%%%%%%%%%%%%%%%%%%%%%%%%%%%%%%%%%%%%
\frontmatter

% Puedes agregar dedicatorias desde el YAML si quieres, o dejarlas fijas aquí

%%%%%%%%%%%%%%%%%%%%%%%%%%%%%%%%%%%%%%%%%%%%%%%%%%%%%
%                   ÍNDICES                         %
%%%%%%%%%%%%%%%%%%%%%%%%%%%%%%%%%%%%%%%%%%%%%%%%%%%%%
\setcounter{secnumdepth}{3}
\setcounter{tocdepth}{3}

\tableofcontents



\cleardoublepage

  \cleardoublepage       % Empieza en página derecha (estándar en tesis)
  \chapter*{Resumen}     % Crea el título "Resumen" sin numeración (Capítulo *)
  \addcontentsline{toc}{chapter}{Resumen} % (Opcional) Lo añade al índice
  
  Resumen del trabajo de
investigación\ldots{}             % Aquí se pega el texto que escribiste en el YAML
  
  \cleardoublepage

%%%%%%%%%%%%%%%%%%%%%%%%%%%%%%%%%%%%%%%%%%%%%%%%%%%%%
%                   CONTENIDO (El Corazón)          %
%%%%%%%%%%%%%%%%%%%%%%%%%%%%%%%%%%%%%%%%%%%%%%%%%%%%%
\mainmatter
\def\baselinestretch{1.5}

% --- CAMBIO CRÍTICO 2: Aquí RMarkdown inyecta todo tu texto ---
\chapter{Introducción}

La Fórmula 1 es la máxima categoría del automovilismo mundial, donde la
diferencia entre la victoria y la derrota suele medirse en milésimas de
segundo. Si bien el rendimiento del vehículo y la habilidad del piloto
son fundamentales, el trabajo en equipo durante las paradas en boxes
(\emph{pit stops}) representa un factor crítico en la estrategia de
carrera.

Entre los años 1994 y 2010, la Fórmula 1 atravesó una de sus épocas más
dinámicas, caracterizada por la reintroducción y posterior prohibición
del repostaje de combustible durante las carreras.

Este trabajo de investigación tiene como propósito realizar un análisis
estadístico exhaustivo sobre el conjunto de datos de paradas en boxes de
la F1 en dicho periodo.

A través del uso de herramientas computacionales en R, se busca
descomponer la varianza de los tiempos de parada en función de factores
como la escudería, el año y las estrategias de carrera.

\chapter{Planteamiento del Problema}

En la Fórmula 1, el tiempo de permanencia en el \emph{pit lane} es una
variable de alto impacto que condiciona el resultado deportivo.

Durante el ciclo 1994-2010, las reglas respecto a las operaciones
permitidas en los boxes cambiaron drásticamente, siendo el más crítico
la reintroducción y posterior prohibición del repostaje de combustible.

El problema central de esta investigación consiste en evaluar, desde una
perspectiva estadística, cómo evolucionaron los tiempos de estas paradas
a lo largo de los años y determinar si existen diferencias
significativas en la eficiencia operativa de las distintas escuderías.

La ausencia de un modelo claro que cuantifique cómo estas variables
categóricas impactan en el tiempo esperado de parada genera la necesidad
de responder a las siguientes preguntas de investigación:

\textbf{1.} ¿Cuál fue el tiempo promedio de las paradas en boxes durante
el periodo de adaptación a las nuevas reglas (1994-1996)?

\textbf{2.} ¿Cómo varían los tiempos de los pit stops al compararlos a
lo largo de los distintos años de competencia estudiados?

\textbf{3.} ¿Cuáles son las estrategias de paradas (1, 2 o 3 paradas)
más predominantes a lo largo del periodo evaluado?

\section{Justificación}

El análisis temporal de las paradas en boxes trasciende el ámbito
puramente deportivo; representa un caso de estudio ideal sobre cómo los
equipos de alto rendimiento optimizan procesos bajo condiciones de
extrema presión y variabilidad. Se justifica esta investigación en la
oportunidad de aplicar modelos estadísticos para cuantificar el impacto
real que tienen los cambios de normativas externas (como la prohibición
del repostaje) sobre los tiempos de ejecución mecánica. Desde la
perspectiva académica, este trabajo permite demostrar la utilidad del
lenguaje R para depurar, transformar y extraer conclusiones
significativas a partir de registros cronométricos crudos, consolidando
habilidades prácticas en el manejo de distribuciones asimétricas y el
análisis de varianza.

\section{Objetivos}

\subsection{Objetivo General}
\begin{itemize}
  \item{Analizar los patrones históricos y los factores determinantes en los tiempos de las paradas en boxes de la Fórmula 1 (1994-2010) para caracterizar el comportamiento de la eficiencia mecánica de las escuderías.}
\end{itemize}

\subsection{Objetivos Específicos}
\begin{itemize}
  \item{1. Comparar los tiempos registrados según el orden de la parada (primera parada vs. segunda parada) para analizar las razones estadísticas por las cuales el segundo pit stop tiende a presentar un menor tiempo. }
  \item{2.Describir la influencia de la variable escudería en la distribución de los tiempos de los pit stops. }
  \item{3. Analizar la distribución de frecuencias del número de paradas por piloto para identificar las estrategias de carrera.}
\end{itemize}

\section{Cobertura}

\subsection{Horizontal}

La cobertura horizontal abarca las métricas operativas y características
intrínsecas de cada parada en boxes registrada.

\begin{table}[ht]
\centering
\caption{Variables consideradas} 
\resizebox{17cm}{!} {
\begin{tabular}{ccc}
  \hline
\ \ \ \ \ TIPO&VARIABLE&DESCRIPCIÓN \\ 
\hline
          Categórica & Car & Escudería o equipo de mecánicos (e.g., Ferrari, McLaren).\\ 
          Numérica & Year & Año en que se realizó el campeonato.\\ 
          Numérica & Stops & Número de parada en esa carrera para ese piloto.\\ 
          Numérica & Time & Tiempo exacto que transcurre desde que el vehículo entra al pit lane hasta que sale.\\ 
  \hline
\end{tabular}
}
\end{table}

\subsection{Vertical}

La cobertura vertical comprende un total de 9.921 observaciones. Cada
fila representa una parada en boxes única ejecutada durante un Gran
Premio de Fórmula 1. El alcance temporal abarca desde el año 1994 hasta
el año 2010.

\section{Periodo de Referencia}

La fuente de los datos corresponde a un registro histórico estructurado.
El periodo de referencia efectivo para el análisis descriptivo se centra
en el intervalo 1994-2010.

\section{Variables}

La variable principal en estudio es el Tiempo de Parada (Time), la cual
representa la eficiencia operativa del equipo de mecánicos, medida en
los segundos transcurridos durante cada pit stop.

\chapter{Marco Teórico}
\subsection{Bases Teóricas}

El presente análisis se fundamenta en las técnicas de la Estadística
Descriptiva, la cual permite recolectar, organizar, resumir y presentar
los datos de manera informativa. Los principales estadísticos y gráficos
empleados son:

\subsection{Medidas de Tendencia Central)}

Se utilizará la Media Aritmética (promedio de los tiempos de parada) y
la Mediana (el valor central que divide la distribución en dos partes
iguales, siendo robusta ante valores atípicos o outliers).

\[ \bar{x} = \frac{1}{n} \sum_{i=1}^{n} x_i \] \[ Me = \begin{cases} 
x_{\frac{n+1}{2}} & \text{si } n \text{ es impar} \\ 
\frac{x_{\frac{n}{2}} + x_{\frac{n}{2} + 1}}{2} & \text{si } n \text{ es par} 
\end{cases} \]

\subsection{Medidas de Dispersión}

Se aplicará la Desviación Estándar para cuantificar el grado de
variabilidad o dispersión de los tiempos de parada en los boxes respecto
a la media.

\[ s = \sqrt{\frac{1}{n-1} \sum_{i=1}^{n} (x_i - \bar{x})^2} \]

\subsection{Distribución de Frecuencias}

Tablas que resumen cuántas veces ocurre un evento (frecuencia absoluta)
y su proporción (frecuencia relativa), vitales para analizar el número
de paradas por piloto.

\[ f_r = \frac{f_i}{n} \]

\subsection{Gráficos}

Gráfico de Cajas (Boxplot): Herramienta visual que representa la
distribución de los datos a través de sus cuartiles. Es fundamental en
este estudio para comparar visualmente los tiempos de la primera y
segunda parada, identificando rápidamente la mediana y los valores
atípicos.

Gráfico de Barras: Representación gráfica para variables categóricas o
discretas (como Escuderías o número de paradas) que muestra la
frecuencia o promedios mediante barras rectangulares.

\chapter{Marco Metódico}

En esta sección se detalla el proceso de Extracción, Transformación y
Carga (ETL) aplicado al conjunto de datos crudos de la Fórmula 1. La
preparación adecuada de los datos es un paso ineludible antes de aplicar
los modelos descritos en el marco teórico.

\section{Bases metódicas utilizadas en el trabajo de investigación}

El procesamiento de la base de datos se realizó utilizando el entorno R
y el ecosistema de paquetes \texttt{tidyverse}. A continuación, se
detalla el proceso de Extracción, Transformación y Carga (ETL) aplicado
al conjunto de datos crudos de la Fórmula 1.\\

\subsection{1. Carga de Datos}

El primer paso consistió en la importación del archivo plano que
contiene los registros históricos de las paradas en boxes a lo largo de
los diferentes Grandes Premios.

\begin{Shaded}
\begin{Highlighting}[]
\CommentTok{\# Importación del dataset original}

\NormalTok{f1\_datos }\OtherTok{\textless{}{-}} \FunctionTok{read\_csv}\NormalTok{(}\StringTok{"dataset/formula1.csv"}\NormalTok{)}
\end{Highlighting}
\end{Shaded}

\subsection{Limpieza y Tratamiento de Valores Nulos (NAs)}

\begin{Shaded}
\begin{Highlighting}[]
\CommentTok{\# Eliminación de valores nulos en la variable objetivo}
\NormalTok{f1\_limpio }\OtherTok{\textless{}{-}}\NormalTok{ f1\_datos }\SpecialCharTok{\%\textgreater{}\%}
  \FunctionTok{drop\_na}\NormalTok{(Time)}
\end{Highlighting}
\end{Shaded}

Transformación de Tipos de Datos y Filtrado

\subsubsection{Item}

\subsection{Item n}

Descripción\\

\chapter{Análisis de Resultados}

\{\small Este capítulo presenta los resultados obtenidos al aplicar las
metodologías descritas en el capítulo anterior, así como del análisis
estadístico descriptivo de las variables estudiadas y las ventajas y
desventajas al aplicar dichos tests.\}

\section{Presentación de resultados 1}

Descripción\\

\section{Presentación de resultados 2}

Descripción\\

\section{Presentación de resultados n}

Descripción\\

\chapter{Conclusiones y Recomendaciones}

\textbf{1.-} Conclusion 1 del trabajo de investigación.\\

\textbf{Recomendación} Recomendación 1 del trabajo de investigación.\\

\textbf{2.-} Conclusion 2 del trabajo de investigación.\\

\textbf{Recomendación} Recomendación 2 del trabajo de investigación.\\

\textbf{n.-} Conclusion n del trabajo de investigación.\\

\textbf{Recomendación} Recomendación n del trabajo de investigación.\\

\appendix
\renewcommand{\thechapter}{\Alph{chapter}}

\chapter{Anexo A}

\chapter{Anexo B}

\chapter{Anexo C}

%%%%%%%%%%%%%%%%%%%%%%%%%%%%%%%%%%%%%%%%%%%%%%%%%%%%%
%                   REFERENCIAS                     %
%%%%%%%%%%%%%%%%%%%%%%%%%%%%%%%%%%%%%%%%%%%%%%%%%%%%%
% Mantenemos la estructura natbib de tu plantilla
\renewcommand{\bibname}{Lista de Referencias}
\bibliographystyle{apalike} 
\bibliography{bibliografia.bib} 

%%%%%%%%%%%%%%%%%%%%%%%%%%%%%%%%%%%%%%%%%%%%%%%%%%%%%
%                   APÉNDICES                       %
%%%%%%%%%%%%%%%%%%%%%%%%%%%%%%%%%%%%%%%%%%%%%%%%%%%%%

\end{document}
